%!TEX root = ../main/main.tex

\subsection{Glushkov Construction with Variable Captures}
\label{subsec:glushkov-captures}

We formulate a construction that applies the position-automaton
construction of Section~\ref{subsec:glushkov} to an expression over
letters and variable markers, and then interprets the resulting
transitions as VA transitions. Throughout, we use the document positions
and validity conditions defined in Section~\ref{sec:preliminaries}.

\paragraph{1. Input expressions}
Let $\gamma$ be a regex formula over $\Sigma$ and $\mathcal V$, and let
$E=\tau(\gamma)$ be its translation into an ordinary regular expression
over $\Omega=\Sigma\cup\Gamma_{\mathcal V}$, as defined in
Section~\ref{subsec:tau}. We use the partial-mapping semantics of
Section~\ref{subsec:mappings} and allow empty captures. The expression
is matched against the entire document. The construction below applies
the position automaton of Section~\ref{subsec:glushkov} to $E$ and then
interprets its transitions as VA transitions. If the input is already
written with explicit markers, we take that expression directly as $E$;
otherwise we set $E=\tau(\gamma)$. We do not impose the well-designedness
restriction of REQL here; the construction is defined for arbitrary regex
formulas. That restriction becomes relevant when sequentiality and the
logical-VA model are discussed.

\paragraph{2. Positions of letters and markers}
Number every occurrence of a symbol of $\Omega$ in $E$ from left to right.
If there are $n$ occurrences, let $\operatorname{Pos}(E)=\{1,\ldots,n\}$,
and let $\ell_E(i)\in\Omega$ be the label at position $i$.
Replacing every symbol occurrence by its position produces $\widehat{E}$.
Different occurrences receive different positions even when their labels
coincide. Neither $\varnothing$ nor $\varepsilon$ receives a position.
For example, $x\vdash\cdot a\cdot b\cdot\dashv x$ has four positions,
labelled $x\vdash$, $a$, $b$, and $\dashv x$, respectively.

\paragraph{3. First, last, and follow}
Use the language-based definitions of $\operatorname{first}$,
$\operatorname{last}$, and $\operatorname{follow}$ from
Section~\ref{subsec:glushkov}, applied to the marked expression
$\widehat{E}$ over $\operatorname{Pos}(E)$. Thus, the construction is
specified in terms of which position sequences occur in
$\mathcal L(\widehat{E})$. At this stage, each variable marker is treated
as an ordinary symbol at its position. Its zero-width interpretation as a
VA transition is used only in Step~5.

\paragraph{4. Constructing the position automaton}
Apply the Glushkov construction over $\Omega$ using the sets above.
The states are $Q_E=\{0\}\cup\operatorname{Pos}(E)$, with initial state $0$.
Define
\[
\begin{aligned}
 \delta_E={}&\{(0,\ell_E(j),j)\mid j\in\operatorname{first}(E)\}\\
            &{}\cup\{(i,\ell_E(j),j)\mid
                i\in\operatorname{Pos}(E),\ j\in\operatorname{follow}(i,E)\},
\end{aligned}
\]
and
\[
 F_E=
 \begin{cases}
   \operatorname{last}(E)\cup\{0\},&\text{if }\varepsilon \in \mathcal L(E),\\
   \operatorname{last}(E),&\text{otherwise}.
 \end{cases}
\]
Marking preserves nullability, so
$\varepsilon\in\mathcal L(E)$ if and only if
$\varepsilon\in\mathcal L(\widehat{E})$.
At this stage, the position automaton reads words over $\Omega$: each
transition reads one letter or one marker as an ordinary alphabet symbol.
There is one state per symbol occurrence and one additional initial state.

\paragraph{5. Interpreting the automaton as a VA}
For a regex formula $\gamma$ with $E=\tau(\gamma)$, define the
associated \emph{Glushkov VA} to be
\[
 G_\gamma=(Q_E,0,F_E,\delta_E).
\]
For an input already written with explicit markers, use the same tuple
and denote it by $G_E$.
The transition relation is now interpreted according to
Definition~\ref{def:configuration}:
\begin{enumerate}
 \item a transition labelled by $a\in\Sigma$ reads the next document
       character, requires it to equal $a$, and advances one position;
 \item a transition labelled by $x\vdash$ opens $x$ at the current
       document position without advancing;
 \item a transition labelled by $\dashv x$ closes $x$ at the current
       document position without advancing.
\end{enumerate}
Runs start at document position $1$ and finish at position $|d|+1$.
A run is accepting when its last state belongs to $F_E$.
Its validity and induced mapping are those of
Definitions~\ref{def:valid-run} and~\ref{def:run-mapping}. Accordingly,
\[
 \llbracket G_\gamma\rrbracket_d
 =\{\mu_\rho\mid\rho\text{ is a valid accepting run of }G_\gamma
                       \text{ over }d\}.
\]
The same convention applies to $G_E$. Opening and closing a variable at
the same document position produces an empty span. In particular,
$x\{\varepsilon\}$ translates to a nonempty word of two markers, yet
its VA can accept the empty document and capture $[1,1\rangle$.
Variable transitions retain their marker labels; no unlabelled
$\varepsilon$-transitions are introduced.

\paragraph{6. Variables and induced mappings}
The variables of $G_\gamma$ are exactly those of $\gamma$, that is,
$\operatorname{var}(G_\gamma)=\operatorname{var}(\gamma)$. Every
occurrence of a capture $x\{\gamma_0\}$ contributes one opening marker
position and one closing marker position in $E=\tau(\gamma)$; these
positions are states of $G_\gamma$. A valid run that traverses the
opening position at document position $i$ and the closing position at
document position $j$ induces the span $[i,j\rangle$ for $x$, as in
Definition~\ref{def:run-mapping}. Thus the mappings are read directly
from the construction; no separate annotation of captures is needed.

\begin{example}[Glushkov construction with a variable capture]
\label{ex:glushkov-capture}
Let $\Sigma=\{a,b\}$ and consider
\[
 \gamma=x\{(a+b)^*\}\cdot a.
\]
The expression captures a possibly empty prefix in $x$, followed by a final
$a$ outside the capture. Translating the capture gives
\[
 E=\tau(\gamma)=x\vdash\cdot(a+b)^*\cdot\dashv x\cdot a.
\]
There are five symbol occurrences, with labels
\[
\begin{array}{c|ccccc}
 i & 1 & 2 & 3 & 4 & 5\\ \hline
 \ell_E(i) & x\vdash & a & b & \dashv x & a
\end{array}
\]
The corresponding marked expression is
\[
 \widehat{E}=1\cdot(2+3)^*\cdot4\cdot5.
\]
Positions $2$ and $5$ are distinct occurrences of $a$. The subexpression
$(2+3)^*$ is nullable, but $\widehat{E}$ is not: positions $1$, $4$, and $5$ must
still be read. Consequently,
\[
 \varepsilon \notin \mathcal L(\widehat{E}),\qquad
 \operatorname{first}(E)=\{1\},\qquad
 \operatorname{last}(E)=\{5\}.
\]
The follow sets are
\[
\begin{array}{c|c}
 i & \operatorname{follow}(i,E)\\ \hline
 1 & \{2,3,4\}\\
 2 & \{2,3,4\}\\
 3 & \{2,3,4\}\\
 4 & \{5\}\\
 5 & \varnothing
\end{array}
\]
After opening $x$, the run may read $a$ or $b$, or close $x$ immediately.
After a symbol inside the star, it may read another symbol of the star or
close the capture. Closing $x$ must be followed by the final $a$.

The resulting VA is $G_\gamma=(Q_E,0,F_E,\delta_E)$, with
\[
 Q_E=\{0,1,2,3,4,5\},\qquad F_E=\{5\}.
\]
Its complete transition relation is
\[
\begin{aligned}
 \delta_E={}&\{(0,x\vdash,1),(4,a,5)\}\\
 &{}\cup\{(i,a,2),(i,b,3),(i,\dashv x,4)\mid i\in\{1,2,3\}\}.
\end{aligned}
\]
Thus, it has six states and eleven transitions. In particular, the
transition $(1,\dashv x,4)$ is present because the star may be skipped;
it closes the variable without consuming a document character.

On the document $d=aba$, a valid accepting run is
\[
\begin{aligned}
 \rho:\quad
 (0,1)&\xrightarrow{x\vdash}_d(1,1)
       \xrightarrow{a}_d(2,2)
       \xrightarrow{b}_d(3,3)\\
      &\xrightarrow{\dashv x}_d(4,3)
       \xrightarrow{a}_d(5,4).
\end{aligned}
\]
Here, the first component of each configuration is an automaton state and
the second is a document position. The variable opens at position $1$ and
closes at position $3$, yielding
\[
 \mu_\rho=[x\mapsto[1,3\rangle],\qquad d(1,3)=ab.
\]
The last $a$ is consumed after the capture closes. This is the only valid
accepting run on $aba$, so
\[
 \llbracket G_\gamma\rrbracket_{aba}=\{[x\mapsto[1,3\rangle]\}.
\]
The same automaton also illustrates an empty capture. On $d=a$, the run
\[
 (0,1)\xrightarrow{x\vdash}_d(1,1)
       \xrightarrow{\dashv x}_d(4,1)
       \xrightarrow{a}_d(5,2)
\]
is valid and accepting and induces $[x\mapsto[1,1\rangle]$.
Both markers are traversed even though the captured substring is empty.
\end{example}

\begin{theorem}[Correctness of the Glushkov construction]
\label{theo:glushkov-correct}
For every regex formula $\gamma$ and every document $d$,
\[
  \llbracket G_\gamma\rrbracket_d = \llbracket\gamma\rrbracket_d.
\]
In particular, the underlying NFA of $G_\gamma$ recognises
$\mathcal L(\tau(\gamma))$.
\end{theorem}
The proof combines the equivalence between the two-layer semantics of
$\gamma$ and the word semantics of $\tau(\gamma)$ with the correctness
of the Glushkov construction for ordinary expressions; it is given in
Appendix~\ref{app:glushkov-correctness}.

%comentario6
%agregar nocion de 1 unambiguous con determinista
%que pasa con la version de 1-unambiguous